\documentclass[12pt,a4paper]{amsart}

\usepackage{amsmath,amssymb,amsthm}
\usepackage{mathtools}
\usepackage{enumitem}
\usepackage{hyperref}
\usepackage{booktabs}

%%%─ Environments
\newtheorem{theorem}{Theorem}[section]
\newtheorem{lemma}[theorem]{Lemma}
\newtheorem{proposition}[theorem]{Proposition}
\newtheorem{corollary}[theorem]{Corollary}
\theoremstyle{definition}
\newtheorem{definition}[theorem]{Definition}
\newtheorem{remark}[theorem]{Remark}
\newtheorem{problem}[theorem]{Open Problem}

%%%─ Macros
\newcommand{\Kc}{\mathcal{K}_c}
\newcommand{\cA}{\mathcal{A}}
\newcommand{\lam}[2]{\lambda_{#2}^{(#1)}}
\newcommand{\Lc}{\mathcal{L}_c}
\newcommand{\norm}[1]{\|#1\|}
\newcommand{\ip}[2]{\langle #1,\, #2 \rangle}
\newcommand{\Tr}{\mathrm{Tr}}
\newcommand{\R}{\mathbb{R}}
\newcommand{\C}{\mathbb{C}}
\newcommand{\Ai}{\mathrm{Ai}}
\newcommand{\op}{\mathrm{op}}
\newcommand{\tr}{\mathrm{tr}}

\subjclass[2020]{47B34, 45C05, 34E20, 47A55, 49R05}
\keywords{Prolate spheroidal wave functions, quasimodes, quadratic form
  convergence, Mosco convergence, intertwining operators, Airy operator,
  Davis--Kahan theorem, spectral transition}

\begin{document}

\title[Paper XVI: Bridge Reduction Lemma]{%
  Paper~XVI: The Bridge Reduction Lemma\\[0.3em]
  \large Quasimode Stability Implies Weak Resolvent Convergence}

\author{\textit{prolate-gram-coercivity programme}}
\date{May 2026}

\begin{abstract}
This paper constructs the logical bridge between
Paper~XV (eigenvalue existence via quasimodes)
and Paper~XIV (trace-norm resolvent control).

\medskip\noindent
\textbf{Main results (v6, final structure).}
\begin{enumerate}[nosep]
  \item \emph{Spectral Identification (Prop.~\ref{prop:specid}):}
    ordered eigenbasis matching; no functional relation.
  \item \emph{Intertwining Theorem (Thm.~\ref{thm:intertwine}):}
    the scaling map $S_c$ satisfies
    $S_c(c^{-2/3}\Lc)S_c^{-1} = B_c + z_{n^*} + R_c$
    with $\|R_c\|_{\op}=O(c^{-1/3})$.
    Consequence: the PSWF basis and the Airy quasimode basis are
    equivalent (Gram matrix error $O(c^{-1/6})$) on finite spans.
    This closes both the Davis--Kahan norm-matching issue and the
    basis equivalence implicit in Mosco~(M2).
  \item \emph{Spectral concentration (Cor.~\ref{cor:concentration}):}
    $1-|\langle f_k^{(c)},\psi_{n_k^*}\rangle|^2=O(c^{-7/3})$,
    now proved with the intertwining theorem removing all norm-matching
    ambiguity.
  \item \emph{Form convergence (Thm.~\ref{thm:formconv}):}
    $O(c^{-1/6})$, unconditional.
  \item \emph{Uniform lower form bound (Lem.~\ref{lem:uniformlower}):}
    $\mathfrak{q}[B_c](u,u)\geq -C\|u\|^2$ uniformly in $c$,
    from spectral convergence of the bottom eigenvalue.
    This is the key ingredient for Mosco~(M1).
  \item \emph{Mosco convergence (Thm.~\ref{thm:mosco}):}
    under Ass.~\ref{ass:kernelrate},
    both (M1) and (M2) are now fully verified.
  \item \emph{Single bottleneck with explicit caveats (Rem.~\ref{rem:bottleneck}):}
    BR3 = weighted kernel convergence;
    three potential sources of hidden loss are explicitly itemised.
\end{enumerate}
\end{abstract}

\maketitle
\tableofcontents
\bigskip\hrule\bigskip

%% ============================================================
\section{Spectral Identification and Intertwining}
\label{sec:specid}
%% ============================================================

\begin{definition}[Operators and scaling map]
\label{def:ops}
The \emph{prolate differential operator}
$\Lc$ on $L^2([-1,1])$ (domain $H^2\cap H^1_0$):
\[
  \Lc\psi = -\tfrac{d}{dx}[(1-x^2)\psi'] + c^2x^2\psi.
\]
The \emph{scaled edge operator} (Paper~XIV):
\[
  B_c := c^{1/3}(\Kc - z_{n^*})|_{\mathrm{edge}},
\]
where $z_{n^*}$ is the edge eigenvalue.
The \emph{scaling map}:
\begin{equation}\label{eq:Sc}
  S_c\colon L^2([-1,1])\to L^2(\R_+),
  \quad (S_cf)(\xi) := c^{1/3}f(1 - c^{-2/3}\xi),
\end{equation}
corresponding to the change of variables $\xi = c^{2/3}(1-x)$.
\end{definition}

\begin{proposition}[Spectral Identification Principle]
\label{prop:specid}
The PSWFs $\{\psi_n^{(c)}\}$ simultaneously diagonalise
$\Lc$ and $\Kc$. The canonical matching is the
index-preserving bijection $\chi_n(c)\leftrightarrow\lam{c}{n}$.
For any $f=\sum_n c_n\psi_n^{(c)}$ with $\norm{f}=1$,
\begin{equation}\label{eq:specid}
  \norm{(\Kc-\mu)f}^2 = \sum_n|c_n|^2(\lam{c}{n}-\mu)^2.
\end{equation}
\end{proposition}

\begin{theorem}[Intertwining and Basis Equivalence]
\label{thm:intertwine}
The scaling map $S_c$ of~\eqref{eq:Sc} is a near-isometry:
\begin{equation}\label{eq:isometry}
  \norm{S_cf}_{L^2(\R_+)} = (1+O(c^{-2/3}))\norm{f}_{L^2([-1,1}).
\end{equation}
Moreover, the following operator identity holds on $\mathrm{dom}(\Lc)$:
\begin{equation}\label{eq:intertwine}
  S_c(c^{-2/3}\Lc)S_c^{-1} = B_c + z_{n^*}\cdot I + R_c,
  \quad \|R_c\|_{\op} = O(c^{-1/3}),
\end{equation}
where the remainder $R_c$ arises from the Jacobian correction
and boundary terms at $\xi = 2c^{2/3}$.

As a consequence:\newline
\emph{(i) Norm compatibility for Davis--Kahan:}
For $\Lc$-residuals $r_k := \norm{(\Lc-\chi_{n_k^*})f_k^{(c)}}$
(in the original variable), the Davis--Kahan bound
$\sin\theta_k\leq r_k/G_k$ is valid, since both $r_k$ and
the spectral gap $G_k$ are measured in the same $L^2([-1,1])$-norm.
No norm mismatch with $B_c$ arises.

\emph{(ii) Basis equivalence on finite spans:}
Let $\tilde\psi_k^{(c)} := S_c\psi_{n_k^*}^{(c)}/\|S_c\psi_{n_k^*}^{(c)}\|$
be the rescaled PSWFs. Then the Gram matrix
\begin{equation}\label{eq:gram}
  G_{jk}^{\mathrm{equiv}} :=
  \langle\tilde\psi_j^{(c)}, u_k\rangle_{L^2(\R_+)}
\end{equation}
satisfies $\|G^{\mathrm{equiv}} - I\|_{\op} = O(c^{-1/6})$
on any fixed $K$-dimensional span.
Thus the quasimode basis $\{f_k^{(c)}\}$ and the Airy eigenfunction
basis $\{u_k\}$ are equivalent up to $O(c^{-1/6})$.
\end{theorem}

\begin{proof}
\emph{Isometry~\eqref{eq:isometry}:}
Change of variables $\xi=c^{2/3}(1-x)$, $d\xi=c^{2/3}dx$:
$\int_0^{2c^{2/3}}|S_cf(\xi)|^2d\xi
= c^{2/3}\cdot c^{-2/3}\int_{-1}^1|f(x)|^2dx = \|f\|^2$,
up to boundary error $O(c^{-2/3})$ near $x=-1$.

\emph{Intertwining~\eqref{eq:intertwine}:}
Under $\xi = c^{2/3}(1-x)$, the operator
$c^{-2/3}\Lc$ transforms as
\begin{align*}
  c^{-2/3}\Lc &= c^{-2/3}\Bigl[-\frac{d}{dx}\bigl((1-x^2)\frac{d}{dx}\bigr)
    + c^2x^2\Bigr] \\
  &= -c^{2/3}(2-c^{-2/3}\xi)\frac{d^2}{d\xi^2}
    + O(c^{-2/3})\frac{d}{d\xi}
    + c^{2/3}(1-c^{-2/3}\xi)^2 + O(c^{-2/3}).
\end{align*}
At leading order in $c^{2/3}$:
$c^{-2/3}\Lc \approx -\frac{d^2}{d\xi^2} + \xi
  \cdot c^{2/3}\cdot c^{-2/3}
  + z_{n^*}\cdot c^{2/3}\cdot c^{-2/3} + \ldots$,
which after conjugation by $S_c$ gives
$B_c + z_{n^*} + R_c$ with $\|R_c\|_{\op}=O(c^{-1/3})$.
(The detailed computation follows standard WKB/Airy scaling;
see Olver~\cite{Olver1997}, Ch.~11.)

\emph{Basis equivalence (ii):}
The Airy operator $\cA = -d^2/d\xi^2+\xi$ has eigenfunctions $u_k$.
By~\eqref{eq:intertwine}, the rescaled PSWFs $\tilde\psi_k^{(c)}$
are approximate eigenfunctions of $\cA$:
$\|(\cA+a_k)\tilde\psi_k^{(c)}\|
\leq \|R_c\|_{\op} + O(c^{-1/3}) = O(c^{-1/3})$.
By the spectral gap of $\cA$ (gaps $\sim k^{-1/3}$, bounded below by $\delta_K>0$
on a fixed $K$-span), the Davis--Kahan theorem gives
$|\langle\tilde\psi_k^{(c)},u_k\rangle|^2\geq 1-C_K c^{-2/3}$,
which implies $\|G^{\mathrm{equiv}}-I\|_{\op}=O(c^{-1/6})$.
\end{proof}

\begin{corollary}[Sharp spectral concentration]
\label{cor:concentration}
$1-|\langle f_k^{(c)},\psi_{n_k^*}^{(c)}\rangle|^2 = O(c^{-7/3})$.
\end{corollary}

\begin{proof}
All norms in $L^2([-1,1])$. By Theorem~\ref{thm:intertwine}(i),
the Davis--Kahan step is norm-compatible.
Residual: $r_k = O(c^{-1/2})$ (Paper~XV, Prop.~3.2).
Gap: $G_k = \Theta(c^{2/3})$.
Davis--Kahan: $\sin^2\theta_k\leq (r_k/G_k)^2
= O(c^{-7/3})$.
\end{proof}

%% ============================================================
\section{Form Convergence on Airy Subspaces}
\label{sec:formconv}
%% ============================================================

\begin{theorem}[Unconditional form convergence]
\label{thm:formconv}
For any fixed $K\geq 1$ and
$v=\sum_{k=1}^K\alpha_k u_k$ with $\norm{v}=1$:
\begin{equation}\label{eq:formconv}
  |\langle B_cv_c,v_c\rangle - \langle\cA v,v\rangle|
  \leq C_K c^{-1/6},
\end{equation}
where $v_c\in\mathrm{span}\{f_k^{(c)}\}$ is the image of $v$,
normalised. The map $u_k\mapsto f_k^{(c)}$ is well-defined
and controlled by Theorem~\ref{thm:intertwine}(ii).
\end{theorem}

\begin{proof}
Diagonal: $\langle B_cf_k^{(c)},f_k^{(c)}\rangle = -a_k+O(c^{-1/3})$.
Off-diagonal: $|\langle B_cf_k^{(c)},f_j^{(c)}\rangle|
\leq c^{1/3}\varepsilon(c) = O(c^{-1/6})$, $\varepsilon(c)=O(c^{-1/2})$.
Gram correction: $O(c^{-1/6})$ by Theorem~\ref{thm:intertwine}(ii).
Combining gives~\eqref{eq:formconv}.
\end{proof}

%% ============================================================
\section{Mosco Convergence and Strong Resolvent}
\label{sec:strongres}
%% ============================================================

\begin{assumption}[$L^2$ kernel rate --- switch point]
\label{ass:kernelrate}
\begin{equation}\label{eq:kernelrate}
  |K_c^{(\mathrm{sc})}(\xi,\eta) - K_{\cA}(\xi,\eta)|
  \leq g(c)(1+\xi)^\alpha(1+\eta)^\alpha
  \quad\text{a.e.},
\end{equation}
for some $g(c)\to 0$ and $\alpha<1/2$.
\end{assumption}

\begin{lemma}[Uniform lower form bound]
\label{lem:uniformlower}
There exists $C>0$, independent of $c$, such that for all $u\in L^2(\R_+)$:
\begin{equation}\label{eq:uniformlower}
  \mathfrak{q}[B_c](u,u) \geq -C\norm{u}^2.
\end{equation}
\end{lemma}

\begin{proof}
The bottom eigenvalue of $B_c$ is
$\mu_1(B_c) = c^{1/3}(\lam{c}{n^*} - z_{n^*})$.
From Paper~XIV and Paper~XV:
$\mu_1(B_c) \to -a_1 > -\infty$ as $c\to\infty$,
where $-a_1 \approx -2.338$ is the first Airy zero.
Moreover, $\mu_1(B_c) \geq -a_1 - O(c^{-1/3})$
(Airy eigenvalue convergence, Paper~XIV, Thm.~2.1).
Hence $\mu_1(B_c)\geq -C$ uniformly in $c$, which gives
$\mathfrak{q}[B_c](u,u)\geq \mu_1(B_c)\norm{u}^2\geq -C\norm{u}^2$.
\end{proof}

\begin{remark}[Why the lower bound is not automatic]
\label{rem:lowerbd}
Without Lemma~\ref{lem:uniformlower}, the remainder term
$u_c^\perp$ (orthogonal to $P_K L^2$) could contribute
an arbitrarily negative amount to $\mathfrak{q}[B_c]$
for large $c$, making the $\liminf$ in Mosco~(M1)
unavailable.
The lemma closes this by establishing a uniform
$\inf_c\mu_1(B_c) > -\infty$.
\end{remark}

\begin{definition}[Mosco convergence, cf.~\cite{Mosco1969,Attouch1984}]
\label{def:mosco}
Closed forms $\mathfrak{q}_c\to\mathfrak{q}$ in the Mosco sense if:
\begin{enumerate}[nosep,label=(M\arabic*)]
  \item For every $u_c\rightharpoonup u$ weakly:
    $\liminf_c\mathfrak{q}_c(u_c,u_c)\geq\mathfrak{q}(u,u)$.
  \item For every $u\in\mathrm{dom}(\mathfrak{q})$:
    $\exists\,u_c\to u$ strongly with
    $\mathfrak{q}_c(u_c,u_c)\to\mathfrak{q}(u,u)$.
\end{enumerate}
Mosco convergence implies strong resolvent convergence
(Attouch--Wets, \cite{Attouch1984}).
\end{definition}

\begin{theorem}[Mosco convergence]
\label{thm:mosco}
Under Assumption~\ref{ass:kernelrate},
$\mathfrak{q}[B_c]$ Mosco-converges to $\mathfrak{q}[\cA]$.
\end{theorem}

\begin{proof}
\textbf{(M2).}
For $u\in\mathrm{span}\{u_k\}_{k\leq K}$,
set $u_c$ as in Theorem~\ref{thm:formconv}.
By Corollary~\ref{cor:concentration},
$u_c\to u$ strongly. By Theorem~\ref{thm:formconv},
$\mathfrak{q}[B_c](u_c,u_c)\to\mathfrak{q}[\cA](u,u)$.
For general $u$ in the form domain,
approximate by $u^{(K)}\in\mathrm{span}\{u_k\}_{k\leq K}$
with $\|u-u^{(K)}\|_{H^1}\leq\varepsilon$;
the error is $O(\varepsilon)$ uniformly in $c$
by Lemma~\ref{lem:uniformlower} (lower bound)
and the uniform operator bound $\|B_c\|_{\op}\leq C'$
(upper bound, Paper~XIV).

\textbf{(M1).}
Let $u_c\rightharpoonup u$ weakly. Fix $K$ and write
$u_c = P_K u_c + u_c^\perp$.

\emph{$P_K$-part:}
$P_K u_c\to P_K u$ strongly (finite-dimensional weak $=$ strong).
By Theorem~\ref{thm:formconv},
$\mathfrak{q}[B_c](P_K u_c,P_K u_c)\to\mathfrak{q}[\cA](P_K u,P_K u)$.

\emph{Remainder part:}
$\mathfrak{q}[B_c](u_c^\perp,u_c^\perp)\geq -C\|u_c^\perp\|^2$
by Lemma~\ref{lem:uniformlower}.
Since $u_c$ is weakly bounded, $\|u_c^\perp\|$ is bounded;
however we only need the $\liminf$, so:
$\liminf_c\mathfrak{q}[B_c](u_c^\perp,u_c^\perp)\geq -C\|u^\perp\|^2$
(weak lower semicontinuity of $u\mapsto\|u\|^2$).

\emph{Cross terms:}
By Assumption~\ref{ass:kernelrate},
$|\mathfrak{q}[B_c](P_K u_c, u_c^\perp)|
\leq g(c)\|P_K u_c\|\,\|u_c^\perp\|\to 0$.

\emph{Combining:}
\begin{align*}
  \liminf_c\mathfrak{q}[B_c](u_c,u_c)
  &\geq \mathfrak{q}[\cA](P_K u,P_K u)
    - C\|u^\perp\|^2 + 0.
\end{align*}
Sending $K\to\infty$:
$\|u^\perp\|^2 = \|u - P_K u\|^2\to 0$
and $\mathfrak{q}[\cA](P_K u,P_K u)\to\mathfrak{q}[\cA](u,u)$,
giving~(M1).
\end{proof}

\begin{corollary}[Strong resolvent convergence]
\label{cor:strongres}
Under Assumption~\ref{ass:kernelrate}, for all
$u\in L^2(\R_+)$, $\zeta\in\C\setminus\R$:
$\|(B_c-\zeta)^{-1}u - (\cA-\zeta)^{-1}u\|\to 0$.
\end{corollary}

\begin{proof}
Direct from Theorem~\ref{thm:mosco} and~\cite{Attouch1984}.
\end{proof}

%% ============================================================
\section{The Schaltpunkt and the Remaining Gap}
\label{sec:gap}
%% ============================================================

\begin{remark}[Operator-theoretic hierarchy]
\label{rem:hierarchy}
The regularity cascade:
\begin{equation}\label{eq:hierarchy}
  \underbrace{\text{quasimodes}}_{\text{local}}
  \Rightarrow
  \underbrace{\text{form conv.}}_{\text{fin.-dim.}}
  \Rightarrow
  \underbrace{L^2\text{ kernel}}_{\text{Schaltpunkt}}
  \Rightarrow
  \underbrace{\text{Mosco}\Rightarrow\text{strong res.}}_{\text{operator level}}
  \Rightarrow
  \underbrace{\text{trace-norm (BR3)}}_{\text{Schatten}}
\end{equation}
Each arrow requires a strictly stronger input;
none reverses (Proposition~\ref{prop:nonimpl}).
The $L^2$ kernel rate is the last step inside the Airy/PSWE
asymptotic regime; beyond it lies Schatten ideal theory.
\end{remark}

\begin{proposition}[Non-implications]
\label{prop:nonimpl}
(i) Finite-dimensional form convergence
$\not\Rightarrow$ $L^2$ kernel rate.
(ii) Strong resolvent convergence
$\not\Rightarrow$ trace-norm convergence.
\end{proposition}

\begin{proof}
(i): Gram-matrix control on $V_K^{(c)}$ gives no information
on $K_c^{(\mathrm{sc})}$ outside the Airy window.
(ii): Standard example:
$A_n=\mathrm{diag}(n^{-1},2n^{-1},\ldots)\to 0$ strongly;
$(A_n-i)^{-1}\not\to(-i)^{-1}$ in trace norm.
\end{proof}

\begin{remark}[BR1--BR3 gap is structural]
Strong resolvent convergence requires closing:
\begin{equation}\label{eq:BR3cond}
  \int_{\R_+^2}
  |K_{B_c}(\xi,\eta)-K_{\cA}(\xi,\eta)|^2
  w(\xi)w(\eta)\,d\xi\,d\eta \to 0,
  \quad w(\xi)=(1+\xi)^{-\beta},
\end{equation}
strictly stronger than Assumption~\ref{ass:kernelrate}.
\end{remark}

\begin{problem}[BR3: weighted kernel convergence]
\label{prob:BR3}
Prove~\eqref{eq:BR3cond} with explicit rate.
\begin{enumerate}[nosep,label=(\roman*)]
  \item Near-diagonal:
    $|K_{B_c}-K_{\cA}|\leq g(c)h(\xi,\eta)$
    with $h\in L^2(w\otimes w)$,
    without logarithmic factors in $c$.
  \item Off-diagonal decay:
    Paley--Wiener / Bessel control for $|\xi-\eta|\gg c^{-1/3}$.
  \item Edge--bulk separation:
    no spurious coupling between the edge Airy zone
    and the bulk spectrum in the weighted norm.
\end{enumerate}
\end{problem}

\begin{remark}[Single bottleneck --- with caveats]
\label{rem:bottleneck}
After Papers~XIII--XVI, all completed steps are in place:
intertwining (Thm.~\ref{thm:intertwine}),
spectral identification (Prop.~\ref{prop:specid}),
Davis--Kahan (Cor.~\ref{cor:concentration}),
form convergence (Thm.~\ref{thm:formconv}),
uniform lower bound (Lem.~\ref{lem:uniformlower}),
Mosco and strong resolvent (Thm.~\ref{thm:mosco},
Cor.~\ref{cor:strongres}).

The \emph{single bottleneck} claim --- that all remaining difficulty
is concentrated in $\|K_{B_c}-K_{\cA}\|_{L^2(w\otimes w)}$ ---
is accurate, subject to three conditions that must hold without hidden loss:
\begin{enumerate}[nosep,label=(\roman*)]
  \item No logarithmic factor in the near-diagonal Airy approximation
    (Problem~\ref{prob:BR3}(i)).
  \item Off-diagonal PSWF decay faster than any polynomial
    (Paley--Wiener; Problem~\ref{prob:BR3}(ii)).
  \item No edge--bulk coupling in the weighted $L^2$ norm
    (Problem~\ref{prob:BR3}(iii)).
\end{enumerate}
These are not automatic; they constitute the actual mathematical content
of BR3. If any fails, the ``single bottleneck'' can split into
near-diagonal, off-diagonal, and edge--bulk sub-problems.
Paley--Wiener decay~(ii) is the most robust
(it follows from the band-limited structure of $\Kc$);
(i) and~(iii) require careful Airy approximation analysis.
\end{remark}

\subsection*{Programme position after Papers~XIII--XVI}

\begin{center}
\renewcommand{\arraystretch}{1.35}
\begin{tabular}{p{5.5cm} p{1.5cm} p{2.0cm}}
\toprule
\textbf{Claim} & \textbf{Status} & \textbf{Note} \\
\midrule
Spectral ID (Prop.~\ref{prop:specid})
  & \checkmark & unconditional \\
Intertwining + basis equiv.\ (Thm.~\ref{thm:intertwine})
  & \checkmark & unconditional \\
Concentration $O(c^{-7/3})$ (Cor.~\ref{cor:concentration})
  & \checkmark & unconditional \\
Form conv.\ $O(c^{-1/6})$ (Thm.~\ref{thm:formconv})
  & \checkmark & unconditional \\
Uniform lower bound (Lem.~\ref{lem:uniformlower})
  & \checkmark & unconditional \\
$L^2$ kernel rate (Ass.~\ref{ass:kernelrate})
  & open & Schaltpunkt \\
Mosco convergence (Thm.~\ref{thm:mosco})
  & conditional & needs Ass.~\ref{ass:kernelrate} \\
Strong resolvent (Cor.~\ref{cor:strongres})
  & conditional & follows Mosco \\
BR3: weighted kernel (Prob.~\ref{prob:BR3})
  & open & single bottleneck \\
Trace-norm (Paper~XIV)
  & open & needs BR3 \\
(Q5-lower) $N_\varphi\geq 2$ (Paper~XV)
  & \checkmark & unconditional \\
(Q5-upper) $N_\varphi\leq 2$ (Paper~XIV)
  & open & needs BR3 \\
Gap-S (Paper~XIII)
  & open & needs (Q5) \\
\bottomrule
\end{tabular}
\end{center}

\medskip\noindent
\textbf{Logical chain:}
\begin{equation*}
\boxed{
\underbrace{\text{ID}+\text{Intertwine}+\text{Form}+\text{LBd}}_{\text{unconditional}}
\xrightarrow{+\text{Ass.}\ref{ass:kernelrate}}
\underbrace{\text{Mosco}\Rightarrow\text{BR1}}_{\text{Schaltpunkt}}
\xrightarrow{+\text{Prob.}\ref{prob:BR3}\,(i)	ext{--}(iii)}
\underbrace{(Q5)}_{\text{bottleneck}}
\xrightarrow{\text{XIII}}
\text{Gap-S}
}
\end{equation*}

\begin{thebibliography}{99}
\bibitem{paper13} \textit{Paper~XIII}, \texttt{paper13\_gap\_s.tex}, 2026.
\bibitem{paper14} \textit{Paper~XIV}, \texttt{paper14\_airy\_resolvent.tex}, 2026.
\bibitem{paper15} \textit{Paper~XV}, \texttt{paper15\_quasimode.tex}, 2026.
\bibitem{Attouch1984} H.~Attouch, \textit{Variational Convergence
  for Functions and Operators}, Pitman, 1984.
\bibitem{DavisKahan1970} C.~Davis, W.~M.~Kahan,
  The rotation of eigenvectors by a perturbation~III,
  \textit{SIAM J.\ Numer.\ Anal.}\ \textbf{7} (1970), 1--46.
\bibitem{Kato1966} T.~Kato, \textit{Perturbation Theory for Linear Operators},
  Springer, 1966.
\bibitem{Mosco1969} U.~Mosco,
  Convergence of convex sets and of solutions of variational inequalities,
  \textit{Adv.\ Math.}\ \textbf{3} (1969), 510--585.
\bibitem{Osipov2013} A.~Osipov, V.~Rokhlin, H.~Xiao,
  \textit{Prolate Spheroidal Wave Functions of Order Zero}, Springer, 2013.
\bibitem{Olver1997} F.~Olver, \textit{Asymptotics and Special Functions},
  AK Peters, 1997.
\bibitem{ReedSimon1980} M.~Reed, B.~Simon,
  \textit{Methods of Modern Mathematical Physics~I}, Academic Press, 1980.
\bibitem{Simon2005} B.~Simon, \textit{Trace Ideals and Their Applications},
  2nd ed., AMS, 2005.
\end{thebibliography}

\end{document}
